% Part: normal-modal-logic
% Chapter: completeness
% Section: frame-completeness

\documentclass[../../../include/open-logic-section]{subfiles}

\begin{document}

\olfileid{nml}{com}{fra}

\olsection{Kasampurnan Tumrap Bingkai}

Teorema kasampurnan kanggo \Log{K} bisa digedhekake marang sistem modal
liyane, sawise kita nuduhake yen model kanonik kanggo logika tartamtu
nduweni sipat bingkai kang cocog.

\begin{thm}\ollabel{thm:completeframeprops}
  Yen logika modal normal $\Sigma$ ngemot salah siji !!{formula}s ing
  sisih kiwa \olref{tab:correspondencetable}, mula model kanonik
  kanggo~$\Sigma$ nduweni sipat kang cocog ing sisih tengen.
\end{thm}

\begin{table}[htp]
  \centering
    \begin{tabular}{| l || l |}
      \hline
      \emph{Yen $\Sigma$ ngemot \dots} & \emph{\dots model kanonik
        kanggo $\Sigma$ iku:} \\
      \hline \hline
      \Ax{D}:\;\; $\Box!A \lif \Diamond !A$ & \quad serial; \\
      \hline
      \Ax{T}:\;\; $\Box!A \lif !A$ &\quad  refleksif;\\
      \hline
      \Ax{B}: \;\;$!A \lif \Box\Diamond!A$ &\quad  simetris; \\
      \hline
      \Ax{4}: \;\; $\Box!A \lif \Box\Box!A$ & \quad transitif; \\
      \hline 
      \Ax{5}: \;\; $\Diamond !A \lif \Box\Diamond!A$& \quad Euklidean.\\
      \hline
    \end{tabular}
    \caption{Kasunyatan korespondensi dhasar.}
    \ollabel{tab:correspondencetable}
\end{table}

\begin{proof}
  Kita bahas saben sipat iki kanthi urut.

  Upamane $\Sigma$ ngemot \Ax{D}, lan jupuken $\Delta \in W^\Sigma$;
  kita kudu nuduhake yen ana $\Delta'$ kang $R^\Sigma
  \Delta\Delta'$. Cukup nuduhake yen $\Box^{-1}\Delta$ iku
  $\Sigma$-konsisten, amarga banjur, miturut Lemma Lindenbaum, ana
  himpunan $\Sigma$-konsisten jangkep $\Delta' \supseteq
  \Box^{-1}\Delta$, lan miturut definisi $R^\Sigma$
  \iftag{prvBox}{}{lan \olref[mod]{lem:box-iff-diamond} }kita duwe
  $R^\Sigma \Delta\Delta'$. Dadi, kanggo kontradiksi, upamane
  $\Box^{-1}\Delta$ iku \emph{ora} $\Sigma$-konsisten, yaiku
  $\Box^{-1}\Delta \Proves[\Sigma] \lfalse$. Miturut \olref[mod]{lem:box2},
  $\Delta \Proves[\Sigma] \Box\lfalse$, lan amarga $\Sigma$ ngemot
  \Ax{D}, uga $\Delta \Proves[\Sigma] \Diamond\lfalse$. Nanging $\Sigma$
  iku normal, mula $\Sigma \Proves \lnot\Diamond \lfalse$
  (\olref[prf][nor]{prop:notDiamondBot}), saengga uga $\Delta
  \Proves[\Sigma] \lnot\Diamond \lfalse$, kang nalisir karo
  konsistensi~$\Delta$.

  Saiki upamane $\Sigma$ ngemot \Ax{T}, lan jupuken $\Delta \in
  W^\Sigma$. Kita arep nuduhake $R^\Sigma \Delta\Delta$,
  yaiku,\iftag{prvBox}{}{ miturut \olref[mod]{lem:box-iff-diamond},}
  $\Box^{-1}\Delta \subseteq \Delta$. Nanging yen $\Box!A \in \Delta$,
  miturut \Ax{T} uga $!A \in \Delta$, kaya kang dikarepake.

  Saiki upamane $\Sigma$ ngemot \Ax{B}, lan upamane $R^\Sigma
  \Delta\Delta'$ kanggo $\Delta$, $\Delta' \in W^\Sigma$. Kita kudu
  nuduhake $R^\Sigma \Delta'\Delta$, yaiku,\iftag{prvBox}{
    $\Box^{-1}\Delta' \subseteq \Delta$. Miturut
    \olref[mod]{lem:box-iff-diamond}, iki ekuivalen karo}{}
  $\Diamond\Delta \subseteq \Delta'$. Dadi upamane $!A \in \Delta$.
  Miturut \Ax{B}, uga $\Box\Diamond!A \in \Delta$. Miturut hipotesis
  $R^\Sigma \Delta\Delta'$\iftag{prvBox}{}{ lan
    \olref[mod]{lem:box-iff-diamond}}, kita duwe $\Box^{-1}\Delta
  \subseteq \Delta'$, lan mula $\Diamond!A \in \Delta'$, kaya kang
  dikarepake.

  Saiki upamane $\Sigma$ ngemot \Ax{4}, lan upamane $R^\Sigma
  \Delta_1\Delta_2$ lan $R^\Sigma \Delta_2\Delta_3$. Kita kudu nuduhake
  $R^\Sigma \Delta_1\Delta_3$. Saka hipotesis\iftag{prvBox}{}{ lan
    \olref[mod]{lem:box-iff-diamond}} kita entuk loro-lorone:
  $\Box^{-1}\Delta_1 \subseteq \Delta_2$ lan $\Box^{-1}\Delta_2
  \subseteq \Delta_3$. Kanggo nuduhake $R^\Sigma \Delta_1\Delta_3$,
  cukup nuduhake $\Box^{-1}\Delta_1 \subseteq \Delta_3$. Dadi jupuken
  $!B \in \Box^{-1}\Delta_1$, yaiku $\Box!B \in \Delta_1$. Miturut
  \Ax{4}, uga $\Box\Box!B \in \Delta_1$ lan saka hipotesis kita entuk,
  kawitan, $\Box!B \in \Delta_2$ lan, kapindho, $!B \in \Delta_3$,
  kaya kang dikarepake.

  Saiki upamane $\Sigma$ ngemot \Ax{5}, lan upamane $R^\Sigma
  \Delta_1\Delta_2$ lan $R^\Sigma \Delta_1\Delta_3$. Kita kudu nuduhake
  $R^\Sigma \Delta_2\Delta_3$. Hipotesis pisanan menehi
  $\Box^{-1}\Delta_1 \subseteq \Delta_2$\iftag{prvBox}{}{ miturut
    \olref[mod]{lem:box-iff-diamond}}, lan hipotesis kapindho ekuivalen
  karo $\Diamond\Delta_3 \subseteq \Delta_1$\iftag{prvBox}{,
    miturut \olref[mod]{lem:box-iff-diamond}}{}. Kanggo nuduhake $R^\Sigma
  \Delta_2\Delta_3$\iftag{prvBox}{, miturut
    \olref[mod]{lem:box-iff-diamond}, cukup}{ kita kudu}
  nuduhake $\Diamond\Delta_3 \subseteq \Delta_2$. Dadi jupuken
  $\Diamond!A \in \Diamond\Delta_3$, yaiku $!A \in \Delta_3$. Miturut
  hipotesis kapindho $\Diamond!A \in \Delta_1$ lan miturut \Ax{5}, uga
  $\Box\Diamond!A \in \Delta_1$. Saiki hipotesis pisanan menehi
  $\Diamond!A \in \Delta_2$, kaya kang dikarepake.
  \end{proof}

Minangka akibat, kita entuk asil kasampurnan kanggo sawatara sistem.
Upamane, kita ngerti yen $\Log{S5} = \Log{KT5} =
\Log{KTB4}$ sampurna gegayutan karo kelas kabeh model refleksif lan
Euklidean, kang padha karo kelas kabeh model refleksif, simetris, lan
transitif.

\begin{thm}\ollabel{thm:generaldet}
  Upamane $\mClass{C}_\Ax{D}$, $\mClass{C}_\Ax{T}$,
  $\mClass{C}_\Ax{B}$, $\mClass{C}_\Ax{4}$, lan
  $\mClass{C}_\Ax{5}$ iku, miturut urutane, kelas kabeh model serial,
  refleksif, simetris, transitif, lan Euklidean. Banjur kanggo sembarang
  skema $!A_1$, \dots, $!A_n$ ing antarane \Ax{D},
  \Ax{T}, \Ax{B}, \Ax{4}, lan \Ax{5}, sistem
  $\Log{K}!A_1 \dots !A_n$ ditemtokake dening
  kelas model $\mClass{C} = \mClass{C}_{!A_1} \cap \dots
  \cap \mClass{C}_{!A_n}$.
\end{thm}

\begin{prop}
  Upamane $\Sigma$ sawijining logika modal normal; banjur:
  \begin{enumerate}
  \item\ollabel{prop:anotherfive-a}%
    Yen $\Sigma$ ngemot skema $\Diamond!A \lif \Box
    !A$, mula model kanonik kanggo $\Sigma$ iku fungsional parsial.
  \item Yen $\Sigma$ ngemot skema $\Diamond!A \liff \Box
    !A$, mula model kanonik kanggo $\Sigma$ iku fungsional.
  \item Yen $\Sigma$ ngemot skema $\Box\Box!A \lif \Box
    !A$, mula model kanonik kanggo $\Sigma$ iku padhet lemah.
  \end{enumerate}
(delengen \olref[frd][acc]{tab:anotherfive} kanggo definisi sipat-sipat
  bingkai iki).
\end{prop}

\begin{proof}
  \begin{enumerate}
  \item Upamane $\Sigma$ ngemot skema $\Diamond !A \lif
    \Box !A$. Kanggo nuduhake yen $R^\Sigma$ fungsional parsial,
    kita kudu mbuktekake yen kanggo sembarang $\Delta_1$, $\Delta_2$,
    $\Delta_3 \in W^\Sigma$, yen $R^\Sigma \Delta_1\Delta_2$ lan
    $R^\Sigma \Delta_1\Delta_3$, mula $\Delta_2=\Delta_3$. Amarga
    $R^\Sigma \Delta_1\Delta_2$, kita duwe $\Box^{-1}\Delta_1
    \subseteq \Delta_2$, lan amarga $R^\Sigma \Delta_1\Delta_3$,
    uga $\Box^{-1}\Delta_1 \subseteq \Delta_3$\iftag{prvBox}{}{,
      loro-lorone miturut \olref[mod]{lem:box-iff-diamond}}. Pepadhan
    $\Delta_2=\Delta_3$ bakal nderek yen kita bisa netepake loro
    inklusi $\Delta_2 \subseteq \Delta_3$ lan $\Delta_3 \subseteq
    \Delta_2$. Kanggo inklusi pisanan, jupuken $!A \in \Delta_2$;
    banjur $\Diamond!A \in \Delta_1$, lan miturut skema kasebut lan
    katutupan deduktif $\Delta_1$ uga $\Box!A \in \Delta_1$.
    Miturut hipotesis $R^\Sigma \Delta_1\Delta_3$, mula
    $!A \in \Delta_3$. Inklusi kapindho padha carane.

  \item Iki langsung nderek saka bagean~\olref{prop:anotherfive-a}
    lan bukti keserialan ing \olref{thm:completeframeprops}.

  \item Upamane $\Sigma$ ngemot skema $\Box\Box!A \lif \Box!A$.
    Kanggo nuduhake yen $R^\Sigma$ padhet lemah, jupuken $R^\Sigma
    \Delta_1\Delta_2$. Kita kudu nuduhake yen ana himpunan
    $\Sigma$-konsisten jangkep $\Delta_3$ kang $R^\Sigma
    \Delta_1\Delta_3$ lan $R^\Sigma \Delta_3\Delta_2$. Tetepa:
    \[
    \Gamma = \Box^{-1}\Delta_1 \cup \Diamond\Delta_2.
    \]
    Cukup nuduhake yen $\Gamma$ iku $\Sigma$-konsisten, amarga banjur,
    miturut Lemma Lindenbaum, bisa disambung dadi himpunan
    $\Sigma$-konsisten jangkep~$\Delta_3$ saengga $\Box^{-1}\Delta_1
    \subseteq \Delta_3$ lan $\Diamond\Delta_2 \subseteq \Delta_3$,
    yaiku $R^\Sigma \Delta_1\Delta_3$\iftag{prvBox}{}{ (miturut
      \olref[mod]{lem:box-iff-diamond})} lan $R^\Sigma
    \Delta_3\Delta_2$\iftag{prvBox}{ (miturut
      \olref[mod]{lem:box-iff-diamond})}{}.

    Kanggo kontradiksi, upamane $\Gamma$ ora konsisten. Banjur ana
    !!{formula}s $\Box!A_1$, \dots, $\Box!A_n \in \Delta_1$
    lan $!B_1$, \dots,~$!B_m \in \Delta_2$ saengga \[!A_1, \dots,
    !A_n, \Diamond!B_1, \dots, \Diamond!B_m \Proves[\Sigma]
    \lfalse.\] Amarga $\Diamond (!B_1 \land \dots \land !B_m) \to
    (\Diamond!B_1 \land \dots \land \Diamond!B_m)$ iku !!{derivable}
    ing saben logika modal normal, kita ngandharake runtutan ing ngisor
    iki, kang nalisir karo konsistensi~$\Delta_2$:
    \begin{align*}
      !A_1, \dots, !A_n, & \Diamond!B_1, \dots,
      \Diamond!B_m \Proves[\Sigma] \lfalse \\
      !A_1, \dots,!A_n & \Proves[\Sigma] (\Diamond!B_1 \land
      \dots \land \Diamond!B_m) \lif \lfalse\\
      & \qquad\text{miturut teorema deduksi}\\
      & \qquad\text{\olref[prf][prp]{prop:derivabilityfacts}\olref[prf][prp]{prop:derivabilityfacts-deduction}, lan \Taut} \\
      !A_1, \dots,!A_n & \Proves[\Sigma]
      \Diamond(!B_1 \land
      \dots \land !B_m) \lif \lfalse\\
      & \qquad \text{amarga $\Sigma$ normal} \\
      !A_1, \dots,!A_n & \Proves[\Sigma]
      \lnot\Diamond (!B_1 \land \dots \land !B_m)\\
      & \qquad\text{miturut \PL} \\
      !A_1, \dots,!A_n & \Proves[\Sigma]
      \Box\lnot (!B_1 \land \dots \land !B_m)\\
      & \qquad\text{$\Box\lnot$ kanggo $\lnot\Diamond$} \\
      \Box!A_1, \dots,\Box!A_n
      & \Proves[\Sigma] \Box\Box \lnot (!B_1 \land \dots \land !B_m)\\
      & \qquad\text{miturut \olref[mod]{lem:box1}} \\
      \Box!A_1, \dots,\Box!A_n
      & \Proves[\Sigma]
      \Box\lnot (!B_1 \land \dots \land !B_m)\\
      &\qquad\text{miturut skema $\Box\Box!A \lif \Box!A$} \\
      \Delta_1 & \Proves[\Sigma]
      \Box\lnot (!B_1 \land \dots \land !B_m)\\
      &\qquad\text{miturut monotonisitas, \olref[prf][prp]{prop:derivabilityfacts}\olref[prf][prp]{prop:derivabilityfacts-monotonicity}} \\
      & \Box\lnot (!B_1 \land \dots \land !B_m) \in \Delta_1\\
      &\qquad\text{miturut katutupan deduktif}; \\
      & \lnot (!B_1 \land \dots \land !B_m) \in \Delta_2\\
      &\qquad \text{amarga }
      R^\Sigma \Delta_1\Delta_2. 
    \end{align*}
  \end{enumerate}
\end{proof}

Adhedhasar tuladha-tuladha iki, wong bisa ngira yen saben
sistem~$\Sigma$ logika modal iku \emph{sampurna}, kanthi teges yen
sistem kasebut mbuktekake saben formula kang valid ing saben bingkai
kang ndadekake saben teorema $\Sigma$ valid. Nanging, ana akeh sistem
kang ora sampurna miturut teges iki.

\end{document}
